\section{模型评价与改进}

\subsection{模型优点}

本文模型具有较强的层次性和可解释性。问题一先建立统一功率平衡、绿电指标
和吨氨成本口径，问题二至问题四再依次加入离散开停机、连续负荷率和储能状态
变量，使每一问的变化都能与上一问直接对比。求解方法也与问题结构匹配：
问题二利用开机净成本增量排序，避免枚举大量二元组合；问题三、四采用线性或
混合整数线性模型，便于检查功率平衡、产量约束和购售电互斥约束。所有结果均
保留逐时数据，可追溯到具体场景和时段。

\subsection{模型不足}

模型仍存在一定简化。首先，电解槽和合成氨装置的效率被视为常数，未考虑低负荷
效率变化、启停损耗、爬坡约束和设备寿命折损。其次，24 种风光组合场景按等权
15 天折算全年，不能完全反映真实天气年内相关性。再次，成本中未纳入氨储运、
水耗、检修停机、容量电费和碳交易收益等因素。问题四储能配置以最大弃电场景
为代表，能够体现容量需求上界，但未进一步考虑多日连续阴天和跨日库存。

\subsection{改进方向}

后续可从三方面改进：一是引入更精细的设备运行约束，包括最小开停机时间、
爬坡速率、低负荷效率曲线和制氢--制氨中间库存；二是采用全年逐时风光数据或
概率场景树替代等权日场景，提高年化结果的统计可靠性；三是将绿电直连指标、
辅助服务收益和碳减排收益纳入多目标优化，形成经济性、合规性和系统友好性
之间的 Pareto 权衡方案。
